% !TEX program = xelatex
\documentclass[UTF8,12pt]{ctexart}

\usepackage[a4paper,margin=2.5cm]{geometry}
\usepackage{hyperref}
\usepackage{tabularx}
\usepackage{array}
\usepackage{enumitem}

\hypersetup{
  colorlinks=true,
  linkcolor=black,
  urlcolor=blue
}

\setlist[itemize]{itemsep=0.3em, topsep=0.4em}
\setlist[enumerate]{itemsep=0.3em, topsep=0.4em}
\renewcommand{\arraystretch}{1.2}

\title{autosign.js 小白使用教程（Windows 版）}
\author{}
\date{}

\begin{document}

\maketitle

\noindent 这份教程只讲怎么用，不讲原理。你跟着做完后，可以完成这几件事：

\begin{itemize}
  \item 装好 \texttt{Git} 和 \texttt{Node.js}
  \item 把项目下载到自己电脑
  \item 配好学号和密码
  \item 把签到地点改成自己常用的教学区
  \item 启动 \texttt{autosign.js}
  \item 看懂脚本是不是在正常运行
  \item 需要停掉时知道怎么停
\end{itemize}

\noindent 整个过程都在 Windows 的 PowerShell 里完成。

\section{先准备环境}

\noindent 你电脑里要先有这两个东西：

\begin{itemize}
  \item \texttt{Git}
  \item \texttt{Node.js}
\end{itemize}

\subsection{安装 Git}

\begin{enumerate}
  \item 打开 Git 官网下载页：\url{https://git-scm.com/download/win}
  \item 下载 Windows 安装包。
  \item 一路点下一步安装即可。
  \item 安装完以后，先关闭所有 PowerShell 窗口，再重新打开一个新的 PowerShell。
\end{enumerate}

\subsection{安装 Node.js}

\begin{enumerate}
  \item 打开 Node.js 官网：\url{https://nodejs.org/}
  \item 下载 \texttt{LTS} 版本。
  \item 安装时保持默认选项即可。
  \item 安装完以后，再重新打开一个新的 PowerShell。
\end{enumerate}

\noindent 建议使用 \texttt{Node.js 18+}。

\subsection{检查有没有装好}

\noindent 在 PowerShell 里分别输入下面三条命令：

\begin{verbatim}
git --version
node -v
npm -v
\end{verbatim}

\noindent 如果三条命令都能正常显示版本号，就可以继续。

\section{用 Git 下载项目}

\noindent 因为网络原因，这个项目拉取代码时请直接使用 \url{https://ghfast.top/} 作为 GitHub 代理，不要直接克隆原始 GitHub 地址。

\noindent 先决定你想把项目放在哪个文件夹。下面用 \path|D:\code| 举例，你也可以换成自己的路径。

\begin{verbatim}
cd D:\
mkdir code
cd .\code
git clone https://ghfast.top/https://github.com/5dbwat4/ZJU-live-better.git
cd .\ZJU-live-better
\end{verbatim}

\noindent 执行完后，你现在所在的位置应该就是项目根目录。这个目录里能看到 \texttt{package.json}、\texttt{README.md}、\texttt{courses.zju} 这些内容。

\noindent 以后每次重新打开 PowerShell，想再次运行脚本时，都要先 \texttt{cd} 回这个项目目录。

\noindent 如果你之前已经用原始 GitHub 地址克隆失败过，不用删别的东西，重新执行上面这条带 \texttt{ghfast.top} 的 \texttt{git clone} 就行。

\section{安装依赖}

\noindent 在项目根目录执行：

\begin{verbatim}
npm install
\end{verbatim}

\noindent 第一次安装可能要等几分钟，这很正常。只要命令最终顺利结束，没有停在报错状态，就可以继续。

\section{配置学号和密码}

\noindent 脚本会从项目根目录读取 \texttt{.env} 文件，所以这一步一定要在项目根目录里完成。

\noindent 先打开 \texttt{.env}：

\begin{verbatim}
notepad .env
\end{verbatim}

\noindent 把里面和账号有关的内容改成你自己的：

\begin{verbatim}
ZJU_USERNAME=
ZJU_PASSWORD=
\end{verbatim}

\noindent 注意这几点：

\begin{itemize}
  \item 如果打开后文件里本来就有内容，直接把学号和密码改成你自己的即可。
  \item 这两行里不要加中文引号，也不要在等号前后乱加空格。
  \item 文件名必须叫 \texttt{.env}，不要变成 \texttt{.env.txt}。
  \item \texttt{.env} 要放在项目根目录，也就是和 \texttt{package.json} 同一级的位置。
\end{itemize}

\noindent 改完以后按 \texttt{Ctrl + S} 保存，然后关闭记事本。

\section{地点设置不是必填}

\noindent 这里需要特别说明一下：\texttt{autosign.js} 本身就是自动轮询脚本。

\noindent 对雷达签到来说，它不是“每节课都要你自己选一个地点再签到”。它会在检测到签到后自动尝试点位。现在这套用法里，默认就可以走“自动轮询全部已知点位”的方式，所以你完全可以先不改地点，直接用。

\noindent 也就是说，最省事的用法是：

\begin{itemize}
  \item 配好学号和密码
  \item 直接启动脚本
  \item 让它自己轮询和处理签到
\end{itemize}

\noindent 如果你只是想先用起来，这一节可以直接跳过。

\subsection{可选：设置一个“优先尝试”的地点}

\noindent 如果你希望脚本先优先尝试某个教学区，再继续尝试其他已知点位，可以额外指定一个 \texttt{raderAt}。这不是必填项，只是一个可选优化。

\noindent 当前脚本里已有的地点代号如下：

\begin{center}
\begin{tabularx}{0.9\textwidth}{>{\ttfamily}l X}
\hline
代号 & 地点 \\
\hline
ZJGD1 & 东一教学楼 \\
ZJGX1 & 西教学楼 \\
ZJGB1 & 段永平教学楼 \\
YQ4 & 玉泉教四 \\
YQ1 & 玉泉教一 \\
YQ7 & 玉泉教七 \\
ZJ1 & 之江校区 1 \\
HJC1 & 华家池校区 1 \\
HJC2 & 华家池校区 2 \\
ZJ2 & 之江校区 2 \\
YQSS & 玉泉宿舍点位 \\
ZJG4 & 紫金港大西区 \\
\hline
\end{tabularx}
\end{center}

\noindent 怎么选最简单：

\begin{itemize}
  \item 你平时大多数课在哪个区域，就先填哪个区域。
  \item 如果你最近大多数课都在某一栋教学楼附近，就优先填那一项。
  \item 如果你懒得配，或者上课地点经常变，那就直接用默认的全自动模式。
\end{itemize}

\noindent 现在更推荐用运行命令时传参，而不是手改文件。例如你想让它优先尝试玉泉教一，可以这样启动：

\begin{verbatim}
node courses.zju/autosign.js --raderAt YQ1
\end{verbatim}

\noindent 如果你想完全交给脚本自动轮询全部已知点位，可以显式写成：

\begin{verbatim}
node courses.zju/autosign.js --raderAt AUTO
\end{verbatim}

\section{启动脚本}

\noindent 确认你现在还在项目根目录，然后执行：

\begin{verbatim}
node courses.zju/autosign.js
\end{verbatim}

\noindent 这条命令现在默认就可以按“自动轮询全部已知点位”的方式运行，所以大多数人直接用这一条就够了。

\noindent 启动后注意这几件事：

\begin{itemize}
  \item 这个窗口不要关。
  \item 这个脚本会一直运行，不会自己退出。
  \item 你如果把电脑睡眠了，脚本也会跟着停掉。
  \item 以后每次要重新使用，都要先进入项目目录，再运行这条命令。
\end{itemize}

\section{怎么判断有没有正常运行}

\noindent 这个脚本默认是一直轮询检查，所以只要它持续输出信息，而且没有直接报错退出，通常就说明在正常工作。

\noindent 最常见的正常输出是：

\begin{verbatim}
[Auto Sign-in](Req #1) No rollcalls found.
\end{verbatim}

\noindent 这句话的意思是：

\begin{itemize}
  \item 脚本已经在跑了
  \item 这一次检查没有发现新的签到
\end{itemize}

\noindent 它会持续输出类似的内容，这很正常。

\noindent 如果发现了签到，可能会看到类似下面这些提示：

\begin{verbatim}
[Auto Sign-in](Req #12) Found 1 rollcalls.
[Auto Sign-in] Now answering rollcall #171329
[Auto Sign-in] Detected active rollcall #171329: ...
\end{verbatim}

\noindent 你可以这样理解：

\begin{itemize}
  \item \texttt{Found 1 rollcalls.}：发现有签到了。
  \item \texttt{Now answering rollcall \#...}：脚本已经开始处理这次签到。
  \item \texttt{Detected active rollcall \#...}：确认找到一个正在进行中的签到。
\end{itemize}

\noindent 如果是雷达签到，成功时可能会看到类似：

\begin{verbatim}
[Auto Sign-in] Trying configured Rader location: YQ1 with outcome: ...
\end{verbatim}

\noindent 或者：

\begin{verbatim}
[Auto Sign-in] Congradulations! You are on the call at Rader location: YQ1
\end{verbatim}

\noindent 如果是数字签到，成功时可能会看到：

\begin{verbatim}
[Auto Sign-in] Number Rollcall 171329 succeeded: found code 0123.
\end{verbatim}

\noindent 如果你看到的是下面这种报错：

\begin{verbatim}
[Auto Sign-in](Req #1) Failed to fetch rollcalls:
\end{verbatim}

\noindent 一般说明这次访问课程平台没有成功。先不要慌，优先检查下面几件事：

\begin{itemize}
  \item 当前网络是否正常
  \item 学号和密码有没有填错
  \item 你是不是没有在项目根目录运行脚本
\end{itemize}

\noindent 还有一个容易误会的点：

\begin{itemize}
  \item 默认不开钉钉通知时，终端里不一定会出现“登录成功”字样。
  \item 只要脚本已经开始循环输出 \texttt{No rollcalls found.}，就说明它已经正常跑起来了。
\end{itemize}

\section{怎么停止}

\noindent 如果你想停掉脚本，在当前这个 PowerShell 窗口里按：

\begin{verbatim}
Ctrl + C
\end{verbatim}

\noindent 按完以后，脚本就会停止。

\section{常见问题}

\subsection{输入 \texttt{node -v} 提示找不到命令}

\noindent 说明 \texttt{Node.js} 没装好，或者装完后没有重新打开 PowerShell。

\noindent 你可以这样处理：

\begin{itemize}
  \item 重新安装一次 \texttt{Node.js}
  \item 安装完成后关闭 PowerShell，再开一个新的窗口
  \item 再执行一次 \texttt{node -v}
\end{itemize}

\subsection{输入 \texttt{git --version} 提示找不到命令}

\noindent 说明 \texttt{Git} 没装好，或者装完后没有重新打开 PowerShell。

\noindent 处理方法和上面一样：

\begin{itemize}
  \item 重新安装 \texttt{Git}
  \item 安装完成后关闭 PowerShell，再开一个新的窗口
  \item 再执行一次 \texttt{git --version}
\end{itemize}

\subsection{\texttt{npm install} 失败了}

\noindent 先做这几件事：

\begin{itemize}
  \item 检查网络
  \item 确认你真的在项目根目录
  \item 直接再执行一次 \texttt{npm install}
\end{itemize}

\subsection{明明改了 \texttt{.env}，脚本还是不对}

\noindent 重点检查：

\begin{itemize}
  \item \texttt{.env} 文件是不是放在项目根目录
  \item 文件名是不是 \texttt{.env}
  \item 学号和密码是不是你自己的
  \item 有没有把密码输错
\end{itemize}

\noindent 项目根目录就是能看到 \texttt{package.json} 的那个目录。

\subsection{脚本一直在跑，但签到效果不好}

\noindent 先确认你现在用的是哪种模式：

\begin{itemize}
  \item 如果你是直接运行 \texttt{node courses.zju/autosign.js}，那它现在默认就是自动轮询全部已知点位。
  \item 如果你想让它先优先试某个教学区，可以改成 \texttt{node courses.zju/autosign.js --raderAt YQ1} 这种写法。
\end{itemize}

\noindent 也就是说，\texttt{raderAt} 现在不是必填项，而是一个可选优化项。

\subsection{我下一次怎么快速启动}

\noindent 下次用的时候，最短流程就是：

\begin{verbatim}
cd 你的项目目录
node courses.zju/autosign.js
\end{verbatim}

\noindent 前提是你之前已经做过这三件事：

\begin{itemize}
  \item 装好了 \texttt{Git}
  \item 装好了 \texttt{Node.js}
  \item 在项目里跑过一次 \texttt{npm install}
\end{itemize}

\section{最后再提醒一次}

\noindent 真正最容易漏掉的只有三件事：

\begin{itemize}
  \item \texttt{.env} 里要填你自己的学号和密码
  \item 默认可以直接全自动跑；\texttt{raderAt} 只是可选优化，不是必改项
  \item 启动后终端窗口不要关，停止时按 \texttt{Ctrl + C}
\end{itemize}

\noindent 如果你只是想要一份更短、能直接转发给别人的版本，可以看同目录下的 \texttt{autosign-chat-message.md}。

\end{document}
